% !TeX program = XeLaTeX
% !TeX root = ../pujavidhanam.tex
\chapt{श्री-चित्रगुप्त-पूजा}

% ---------------------------------------------------------------------------
% Source: Tamil-script (Grantha-style) paddhati, "சித்ரகுப்த பூஜை", pp. 6–13.
% Sanskrit reconstructed into Devanagari from the Tamil phonetic spelling.
% The Tamil PROCEDURAL rubrics (madhyāhnika timing; wipe place with cow-dung
% and draw a Chitragupta kolam; keep all items at hand; observe silence until
% the pūjā ends; the upāyana-dāna details below) are NOT typeset here, since the
% house paddhati is Sanskrit-only. They are flagged for the editor's decision.
% Lines marked  % [?]  below are uncertain scan readings — verify against source.
% ---------------------------------------------------------------------------

\input{purvanga/vighneshwara-puja}

\sect{प्रधान-पूजा — श्री-चित्रगुप्त-पूजा}

\shuklambaradharam

\pranayama

% चित्रा-पौर्णमी (Tamil Chithirai / Meṣa, Chitrā-nakṣatra). Verify whether the
% \sankalpa macro emits the नक्षत्र clause; the source specifies चित्रा-नक्षत्र.
\renewcommand{\ayane}{उत्तरायणे}
\renewcommand{\rtu}{वसन्त}
\renewcommand{\masa}{मेष-चैत्र}
\renewcommand{\paksha}{शुक्ल}
\renewcommand{\tithau}{पौर्णम्यां}
\renewcommand{\prityartham}{}
\renewcommand{\additionalSankalpa}{धर्मराज-सभास्थित-श्री-चित्रगुप्त-प्रीत्यर्थं रौरव-महारौरव-कालसूत्र-तामिस्र-अन्धतामिस्र-आदि-नरक-पात-निवृत्त्यर्थं पुण्यलोक-प्राप्त्यर्थं}
\renewcommand{\pujaam}{श्री-चित्रगुप्त-पूजां}

\sankalpa{}

\vighneshvaraYathasthanam

\input{purvanga/aasana-puja}

\input{purvanga/ghanta-puja}

\input{purvanga/kalasha-puja}

\input{purvanga/aatma-puja}

\input{purvanga/pitha-puja}

\input{purvanga/guru-dhyanam}

\sect{षोडशोपचार-पूजा}
\renewcommand{\devAya}{श्री-चित्रगुप्ताय नमः,}
\begin{center}

\twolineshloka*
{चित्रगुप्तं महाप्राज्ञं मध्यस्थं सर्वदेहिनाम्}
{ध्यायामि धर्माधर्मज्ञं लेखिनी-पत्रधारिणम्}
\textbf{श्री-चित्रगुप्तं ध्यायामि।}
\medskip

\twolineshloka*
{आवाहयामि देव त्वाम् अन्तकस्य सभासदम्}
{कायस्थं करुणामूर्तिं कल्मषघ्नं कलाविदम्}
\textbf{श्री-चित्रगुप्तम् आवाहयामि।}
\medskip

\twolineshloka*
{चित्रगुप्त गृहाणेदं चित्रमणि-विनिर्मितम्}
{आसनं हेमरचितं शुभास्तरण-संयुतम्}
\textbf{\devAya{} आसनं समर्पयामि।}
\medskip

\twolineshloka*
{पाद्यं गृहाण परम पापलेश-विवर्जित}
{भाविताशेष-भुवन चित्रगुप्त नमोऽस्तु ते}
\textbf{\devAya{} पाद्यं समर्पयामि।}
\medskip

\twolineshloka*
{अर्घ्यं गृहाण देवेश गन्ध-पुष्पाक्षतैर्युतम्}
{चित्रगुप्त महाप्राज्ञ पापं वारय मेऽखिलम्}
\textbf{\devAya{} अर्घ्यं समर्पयामि।}
\medskip

\twolineshloka*
{गृहाणाचमनं देव करुणारस-सागर}
{त्राहि मां नरकात् घोरात् देहि मे सुकृतां पदम्}
\textbf{\devAya{} आचमनीयं समर्पयामि।}
\medskip

\twolineshloka*
{मधुपर्कं गृहाणेमं महापाप-विनाशन}
{महतामपि पापानां विनाशं कलय प्रभो}
\textbf{\devAya{} मधुपर्कं समर्पयामि।}
\medskip

\twolineshloka*
{पञ्चामृतं गृहाणेदं त्वं चराचर-वल्लभ}
{वञ्चनादि-समुद्भवं पापं मे विनिवारय}
\textbf{\devAya{} पञ्चामृतं समर्पयामि।}
\medskip

\twolineshloka*
{नानानदी-समानीतं नानाधर्म-विचारद}
{स्नानाय गृह्यतां तोयं स्वच्छं गन्धादि-संयुतम्}
\textbf{\devAya{} स्नानं समर्पयामि।}
\medskip

\twolineshloka*
{वस्त्रयुग्मम् इदं देव शस्त्र-हस्त गृहाण मे}
{शास्त्र-लङ्घन-सम्भूतं पापं मे विनिवारय}
\textbf{\devAya{} वस्त्रं समर्पयामि।}
\medskip

\twolineshloka*
{उपवीतम् इदं देव कार्पास-परिकल्पितम्}
{कृतान्त-मित्र गृह्णीष्व कलिजं कलुषं हर}
\textbf{\devAya{} उपवीतं समर्पयामि।}
\medskip

\twolineshloka*
{गन्धं गृहाण सुरभिं उरु-पापविनाशन}
{गतिं देहि परां मह्यं गणय अघानि मा मम}
\textbf{\devAya{} गन्धान् धारयामि।}
\medskip

\twolineshloka*
{अक्षतान् स्वीकुरु विभो यक्ष-गन्धर्व-पूजित}
{अक्षयं देहि मे सौख्यं रक्ष मां करुणानिधे}
\textbf{\devAya{} अक्षतान् समर्पयामि।}
\medskip

\twolineshloka*
{सुविचित्राणि च पुष्पाणि चित्रगुप्त गृहाण भो}
{कमलादीनि रम्याणि विमलं देहि मे पदम्}
\textbf{\devAya{} पुष्पाणि समर्पयामि।}
\medskip

\end{center}

\dnsub{अङ्ग-पूजा}
\begin{longtable}{ll@{— }l}
१.& पापनाशनाय नमः & पादौ पूजयामि\\
२.& कमलवर्णाय नमः & गुल्फौ पूजयामि\\
३.& चित्राय नमः & जानुनी पूजयामि\\
४.& चित्ररूपाय नमः & जङ्घे पूजयामि\\
५.& नताभीष्टदाय नमः & नाभिं पूजयामि\\
६.& वृकोदराय नमः & उदरं पूजयामि\\
७.& प्राज्ञाय नमः & हृदयं पूजयामि\\
८.& कालाय नमः & कण्ठं पूजयामि\\
९.& बहुरूपाय नमः & बाहून् पूजयामि\\
१०.& अन्तकाय नमः & अधरं पूजयामि\\
११.& मध्यस्थाय नमः & मुखं पूजयामि\\
१२.& करुणानिधये नमः & कर्णौ पूजयामि\\
१३.& पद्मनेत्राय नमः & नेत्रे पूजयामि\\
१४.& नानारूपाय नमः & नासिकां पूजयामि\\
१५.& ललिताय नमः & ललाटं पूजयामि\\
१६.& चित्रगुप्ताय नमः & शिरः पूजयामि\\
& सर्वेश्वराय नमः & सर्वाण्यङ्गानि पूजयामि\\
\end{longtable}

\dnsub{नाम-पूजा}
\begin{multicols}{2}
\begin{enumerate}
\item यमाय नमः
\item धर्मराजाय नमः
\item मृत्यवे नमः
\item अन्तकाय नमः
\item वैवस्वताय नमः
\item कालाय नमः
\item सर्वभूतक्षयाय नमः
\item औदुम्बराय नमः
\item दध्नाय नमः
\item नीलाय नमः
\item परमेष्ठिने नमः
\item वृकोदराय नमः
\item चित्राय नमः
\item चित्रगुप्ताय नमः
\end{enumerate}
\end{multicols}
\begin{center}
\textbf{नानाविध-परिमल-पुष्पाणि समर्पयामि।}
\medskip

\twolineshloka*
{धूपं गृहाण धर्मज्ञ गुग्गुलु-संयुतम्}
{आर्तं देहि पदं देव दूरपाप नमोऽस्तु ते}
\textbf{\devAya{} धूपम् आघ्रापयामि।}
\medskip

\twolineshloka*
{साज्यं वर्तित्रयोपेतं शासिताखिल-विष्टप}
{दीपं गृहाण दयया दीनवत्सल ते नमः}
\textbf{\devAya{} दीपं सन्दर्शयामि।}
\medskip

\twolineshloka*
{चित्रान्नं कृसर-पूप-मोदकाज्य-समन्वितम्}
{चित्रगुप्त गृहाणेदं नैवेद्यं निर्मल प्रभो}
\textbf{\devAya{} नैवेद्यं निवेदयामि।}
\medskip

\twolineshloka*
{पूगीफल-समायुक्तं नागवल्ली-दलैर्युतम्}
{कर्पूरचूर्ण-संयुक्तं ताम्बूलं प्रतिगृह्यताम्}
\textbf{\devAya{} ताम्बूलं समर्पयामि।}
\medskip

\twolineshloka*
{नीराजनं गृहाणेदं कर्पूरैः कलितं मया}
{नीतिमार्गचर अनन्त निर्मलाय नमोऽस्तु ते}
\textbf{\devAya{} नीराजनं समर्पयामि।}
\medskip

\twolineshloka*
{प्रकृष्ट-पापनाशाय प्रकृष्ट-फलसिद्धये}
{प्रदक्षिणं करोमि त्वां चित्रगुप्त प्रसीद मे}
\textbf{अनेक-कोटि-प्रदक्षिण-नमस्कारान् समर्पयामि।}
\medskip

\twolineshloka*
{नमस्ते चित्रगुप्ताय धर्मराज-सखाय च}
{लेखिनी-वृत-हस्ताय साक्षिणे सर्वदेहिनाम्}
\textbf{\devAya{} मन्त्रपुष्पं समर्पयामि।}
\medskip

\twolineshloka*
{नानारत्न-समायुक्तं वज्रनाल-समन्वितम्}
{मुक्ताकेसर-संयुक्तं स्वर्णपुष्पं ददाम्यहम्}
\textbf{\devAya{} स्वर्णपुष्पं समर्पयामि।}
\medskip

\twolineshloka*
{यमाय धर्मराजाय मृत्यवे चान्तकाय च}
{वैवस्वताय कालाय सर्वभूतक्षयाय च}

\twolineshloka*
{औदुम्बराय दध्नाय नीलाय परमेष्ठिने}
{वृकोदराय चित्राय चित्रगुप्ताय वै नमो नमः}
\textbf{\devAya{} नमस्कारान् समर्पयामि।}
\medskip

\end{center}

\sect{उत्तराङ्ग-पूजा}
\begin{center}

\twolineshloka*
{चित्रगुप्त कृपाराशे पापं वारय मेऽखिलम्}
{मा मे नरकाः सन्तु पूरयेदं ममेप्सितम्}
\textbf{प्रार्थनां समर्पयामि।}
\medskip

\end{center}

\dnsub{अर्घ्य-प्रदानम्}
अद्य \blank{}-तिथौ श्री-चित्रगुप्त-पूजान्ते क्षीरार्घ्य-प्रदानम् उपायन-दानं च करिष्ये।
\begin{center}

\twolineshloka*
{चित्रगुप्त नमस्तेऽस्तु धर्मराज-सभासद}
{अर्घ्यं गृहाण भगवन् सर्वलोक-शिवङ्कर}
\textbf{श्री-चित्रगुप्ताय नमः, इदम् अर्घ्यम् (त्रिवारम्)।}
\medskip

\twolineshloka*
{वैवस्वत नमस्तुभ्यं धर्माधर्म-विचारक}
{इदम् अर्घ्यं प्रदास्यामि मम पापं विनाशय}
\textbf{वैवस्वताय नमः, इदम् अर्घ्यम् (त्रिवारम्)।}
\medskip

\twolineshloka*
{सर्वलोकाधि-कर्त्रे च सर्वप्राणि-हिताय च}
{तुभ्यम् अर्घ्यं प्रदास्यामि गृहाणेदं चतुर्मुख}
\textbf{ब्रह्मणे नमः, इदम् अर्घ्यम् (त्रिवारम्)।}
\medskip

\end{center}

\dnsub{उपायन-दानम्}
% Tamil rubric (not typeset): in a new winnowing-tray place self-cooked food
% (स्वयम्पाकम्), paper, pencil, tāmbūla and dakṣiṇā, and give to a Brāhmaṇa.
% On that day salt and cow's milk are to be avoided.
चित्रगुप्त-प्रीत्यर्थं स्वयम्पाक-ताम्बूल-दक्षिणादि-सहितम् उपायनं ब्राह्मणाय दद्यात्।

\bigskip
\begin{center}
\textbf{इति श्री-चित्रगुप्त-पूजा समाप्ता।}
\end{center}

\closesection
